\section{Conclusion}

We introduced \benchmark, a comprehensive benchmark for evaluating multimodal embedding models across text, image, video, and visual document modalities. Alongside it, we proposed \model, a strong baseline trained via contrastive learning across a diverse range of tasks and modality combinations. Our extensive experiments demonstrate the effectiveness of \model and the diagnostic value of \benchmark.

% This work opens up several promising future directions. First, we rely on synthetic video data for initial training; future work will incorporate high-quality, real-world video datasets to improve generalization. 
% Second, given the structural similarity between VisDoc and text retrieval tasks, we plan to explore advanced techniques from text embedding literature, including hard negative mining, multi-negative contrastive loss, and query synthesis. 
% In the future, we aim to extend the model and benchmark to include additional modalities and application scenarios, such as audio-text retrieval and surveillance video analysis, to further broaden its impact and utility.

\section{Broader Impact and Limitations} 
Multimodal embedding models are increasingly used in real-world systems such as document retrieval, multimodal search, and content recommendation, where unified representations across heterogeneous modalities enable more accurate and efficient retrieval of relevant content. While our framework advances multimodal embedding learning, the \benchmark benchmark is primarily constructed from English-language data and therefore may not adequately evaluate multilingual or cross-lingual capabilities. In addition, \benchmark focuses on visual modalities and does not cover other important modalities such as audio, limiting its ability to assess truly universal multimodal representations. We view \benchmark as a first step toward comprehensive multimodal evaluation and encourage future work to extend its coverage to additional languages and modalities.
